\section{模型假设}
结合数据组织方式与各模态的物理特性，本文作如下假设，每条假设同时说明其依据。

\begin{enumerate}[label=\arabic*.]
\item 官方转写与视频中的话语顺序一致，转写不改变话语的先后关系。依据是该转写由赛题统一提供并与音轨逐条对应。该假设使词级强制对齐的单调路径约束成立，无法构造字符目标的词仍可按同一顺序单独记录。

\item 同一原词的子词共享父词的时间区间，帧特征在其时间支撑区间内近似为分段常量。依据是词为发音的自然单位，而音频帧率远高于词长：前者使子词继承父词的音视频锚点，后者使式\eqref{q1:eq:pool}的交叠积分退化为加权求和。

\item 在附件2至附件4的对齐版本中，同一序列位置的三个模态共享官方给定的时间对应关系，局部缺失长度以锚点数度量。依据是这些附件按\texttt{aligned\_50}组织，每个序列位置本身即为一个对应单元。问题二与问题三因此只需处理各位置的观测可用性，无需重新估计时间对应关系。

\item 内容位置上的文本\file{[UNK]}与全零音视频向量表示该位置观测不可用，而非有效的特征取值。依据是附件3将局部缺失定义为特征值全为零的连续区间。填充、观测与局部缺失三种状态由此可由掩码唯一确定。

\item 人工遮挡改变可用证据，但不改变样本标签；遮挡后仍以原片段的情感极性与情感强度作为监督目标，以评价输入信息减少时的预测性能。问题二据此构造缺失条件下的训练样本。

\item 标准化统计量仅由附件2的训练集估计，并直接用于验证集与附件3、附件4；推理阶段模型权重与标准化统计量保持固定。依据是各附件由同一特征提取流程生成，特征维度与量纲一致。该约定避免推理阶段引入测试数据信息，并使跨附件指标在同一尺度下可比。
\end{enumerate}